\chapter{并发、线程和同步}

\section{问题入口：两个线程同时计数}

并发最容易制造“偶尔错”的程序。看一个简单计数器：

\begin{lstlisting}[language=C++,caption={数据竞争}]
#include <thread>

int counter = 0;

void increment_many()
{
    for (int i = 0; i < 100000; ++i) {
        ++counter;
    }
}
\end{lstlisting}

如果两个线程同时执行，结果未必是 \code{200000}。\code{++counter} 不是不可分割动作，它通常包含读取、加一、写回。两个线程交错时会丢更新。这是数据竞争，行为不可靠。

\section{mutex 保护共享状态}

\begin{lstlisting}[language=C++,caption={用 mutex 保护计数器}]
#include <mutex>

class Counter {
public:
    void increment()
    {
        std::lock_guard<std::mutex> lock{this->mutex_};
        ++this->value_;
    }

    [[nodiscard]] int value() const
    {
        std::lock_guard<std::mutex> lock{this->mutex_};
        return this->value_;
    }

private:
    mutable std::mutex mutex_;
    int value_ = 0;
};
\end{lstlisting}

\code{std::lock_guard} 是 RAII：构造时加锁，析构时解锁。即使中途抛异常，也会释放锁。

\section{atomic 适合简单共享变量}

简单计数可以用原子类型：

\begin{lstlisting}[language=C++,caption={atomic 计数}]
#include <atomic>

std::atomic<int> counter{0};

void increment_many()
{
    for (int i = 0; i < 100000; ++i) {
        counter.fetch_add(1, std::memory_order_relaxed);
    }
}
\end{lstlisting}

\code{memory_order_relaxed} 只保证计数本身原子，不建立额外同步关系。入门阶段不要随意选择内存序。默认顺序一致最容易推理；只有在测量证明需要时，才进入更细的内存序设计。

\section{条件变量表达等待}

线程间等待不要忙等。生产者消费者可以用 \code{std::condition_variable}。

\begin{lstlisting}[language=C++,caption={条件变量等待队列非空}]
#include <condition_variable>
#include <mutex>
#include <queue>

class WorkQueue {
public:
    void push(int value)
    {
        {
            std::lock_guard<std::mutex> lock{this->mutex_};
            this->queue_.push(value);
        }
        this->ready_.notify_one();
    }

    int pop()
    {
        std::unique_lock<std::mutex> lock{this->mutex_};
        this->ready_.wait(lock, [this] {
            return !this->queue_.empty();
        });
        const int value = this->queue_.front();
        this->queue_.pop();
        return value;
    }

private:
    std::mutex mutex_;
    std::condition_variable ready_;
    std::queue<int> queue_;
};
\end{lstlisting}

\code{wait} 必须带谓词，因为线程可能被虚假唤醒。谓词是等待条件的正式表达。

\section{线程退出也是协议}

生产者消费者队列如果没有退出协议，消费者可能永远阻塞。可以加入关闭标志：

\begin{lstlisting}[language=C++,caption={带关闭协议的队列片段}]
class ClosableQueue {
public:
    void close()
    {
        {
            std::lock_guard<std::mutex> lock{this->mutex_};
            this->closed_ = true;
        }
        this->ready_.notify_all();
    }

    std::optional<int> pop()
    {
        std::unique_lock<std::mutex> lock{this->mutex_};
        this->ready_.wait(lock, [this] {
            return this->closed_ || !this->queue_.empty();
        });
        if (this->queue_.empty()) {
            return std::nullopt;
        }
        const int value = this->queue_.front();
        this->queue_.pop();
        return value;
    }

private:
    std::mutex mutex_;
    std::condition_variable ready_;
    std::queue<int> queue_;
    bool closed_ = false;
};
\end{lstlisting}

关闭不是简单地杀线程，而是让等待方能观察到“不会再有新任务”的状态，并自然退出循环。

\section{本章边界检查}

写并发代码时至少检查：

\begin{itemize}
  \item 哪些状态被多个线程共享？
  \item 每个共享状态由哪个锁或原子保护？
  \item 是否存在锁顺序反转导致死锁？
  \item 条件变量是否用谓词等待？
  \item 线程退出和资源释放路径是否明确？
\end{itemize}

\begin{keyidea}
并发编程的第一原则是减少共享。能在线程内拥有的数据，不共享；必须共享的数据，用清晰同步边界保护。
\end{keyidea}

\section{从一次丢更新推导数据竞争}

\code{++counter} 可以拆成三步：读取旧值、计算新值、写回新值。假设两个线程都看到旧值是 \code{10}：

\begin{lstlisting}[caption={一次丢更新的交错过程}]
线程 A: read counter -> 10
线程 B: read counter -> 10
线程 A: write 11
线程 B: write 11
\end{lstlisting}

执行了两次自增，结果却只增加一次。这不是“线程不够快”，而是没有同步保护共享状态。只要两个线程同时访问同一内存位置，至少一个写，并且没有同步，就形成数据竞争。

\section{锁保护的是不变量，不是一行代码}

给计数器加锁很简单，但真实对象往往有多个成员。锁要保护的是这些成员之间的不变量。比如银行账户转账，不能只给扣款和加款分别加锁后随意组合，还要考虑两个账户的锁顺序。

死锁的典型反例：

\begin{lstlisting}[language=C++,caption={锁顺序反转可能死锁}]
void transfer(Account& from, Account& to, int amount)
{
    std::lock_guard<std::mutex> first{from.mutex};
    std::lock_guard<std::mutex> second{to.mutex};
    from.balance -= amount;
    to.balance += amount;
}
\end{lstlisting}

如果线程一执行 \code{transfer(a, b, 10)}，线程二同时执行 \code{transfer(b, a, 20)}，两个线程可能各自拿住一个锁，然后等待对方释放。可以用 \code{std::scoped_lock} 一次锁住多个互斥量：

\begin{lstlisting}[language=C++,caption={作用域锁同时获取多个锁}]
void transfer(Account& from, Account& to, int amount)
{
    std::scoped_lock lock{from.mutex, to.mutex};
    from.balance -= amount;
    to.balance += amount;
}
\end{lstlisting}

更完整的设计还要处理 \code{from} 和 \code{to} 是同一个账户、余额不足、异常安全等业务规则。并发代码不能脱离不变量讨论。

\section{jthread 让线程生命周期更安全}

\cpp20 提供 \code{std::jthread}，析构时会请求停止并自动 join。相比裸 \code{std::thread}，它更符合 RAII：

\begin{lstlisting}[language=C++,caption={jthread 自动 join}]
#include <chrono>
#include <thread>

void worker(std::stop_token token)
{
    while (!token.stop_requested()) {
        do_one_unit_of_work();
    }
}

void run_background()
{
    std::jthread thread{worker};
    std::this_thread::sleep_for(std::chrono::seconds{1});
} // thread 析构时请求停止并 join
\end{lstlisting}

这不等于可以随便启动后台线程。停止请求只是协作信号，工作函数必须定期检查它。如果线程阻塞在条件变量或 I/O 上，还需要设计唤醒和关闭协议。

\section{条件变量为什么要谓词}

条件变量可能虚假唤醒，也可能在通知前条件已经变化。谓词把“醒来后是否真的能继续”写成可重复检查的条件：

\begin{lstlisting}[language=C++,caption={wait 等价于循环检查}]
while (queue.empty() && !closed) {
    ready.wait(lock);
}
\end{lstlisting}

带谓词的 \code{wait} 帮你写出这个循环。不要用 \code{if} 替代 \code{while}，因为醒来不代表条件一定满足。

\begin{exercisebox}
给 \code{ClosableQueue} 增加一个 \code{push} 规则：关闭后不允许再插入。设计返回值或异常策略，并测试三个场景：正常 push/pop、关闭空队列、关闭后 push。
\end{exercisebox}

\section{并发事故从时间线开始}

并发 bug 的难点是“代码顺序”和“实际交错顺序”不是一回事。两个线程丢更新的时间线可以写得更细：

\begin{longtable}{p{0.12\linewidth}p{0.28\linewidth}p{0.28\linewidth}p{0.14\linewidth}}
\toprule
时刻 & 线程 A & 线程 B & 共享值 \\
\midrule
1 & 读取 counter &  & 10 \\
2 &  & 读取 counter & 10 \\
3 & 计算 11 &  & 10 \\
4 &  & 计算 11 & 10 \\
5 & 写回 11 &  & 11 \\
6 &  & 写回 11 & 11 \\
\bottomrule
\end{longtable}

这个表说明，问题不在循环次数，也不在某个线程“抢占了 CPU”。问题是两个线程同时读写同一个共享状态，且没有同步建立互斥或顺序。任何并发问题都先写时间线：谁读，谁写，哪个状态被共享，哪个同步原语保护它。

\section{锁保护的是状态集合}

只说“给变量加锁”太粗。锁真正保护的是一组状态和它们之间的不变量。以关闭队列为例，共享状态至少有两个：队列内容和关闭标志。

\begin{lstlisting}[language=C++,caption={队列内容和关闭状态必须一起保护}]
std::queue<int> queue_;
bool closed_ = false;
std::mutex mutex_;
\end{lstlisting}

不变量包括：关闭后消费者如果发现队列空，应退出；关闭后是否允许生产者继续 push，要有明确规则；通知必须发生在状态改变之后。于是 \code{closed_} 和 \code{queue_} 应该由同一个互斥量保护。把队列用一个锁、关闭标志用另一个锁，或者把关闭标志改成 atomic 而队列仍用 mutex，都会让推理复杂很多。

\section{push 也需要关闭协议}

前面的 \code{ClosableQueue} 只展示了 \code{close} 和 \code{pop}。完整队列还要定义关闭后 \code{push} 怎么办。这里选择返回 \code{false}：

\begin{lstlisting}[language=C++,caption={关闭后拒绝插入}]
bool push(int value)
{
    {
        std::lock_guard<std::mutex> lock{this->mutex_};
        if (this->closed_) {
            return false;
        }
        this->queue_.push(value);
    }
    this->ready_.notify_one();
    return true;
}
\end{lstlisting}

为什么通知放在释放锁之后？这样等待线程被唤醒后，更可能直接拿到锁继续执行，减少无谓竞争。放在锁内通常也能正确，但要知道自己在做什么。并发代码里，“正确”和“减少等待”是两层问题，先保证正确，再谈性能。

\section{等待谓词是协议文档}

消费者等待的谓词：

\begin{lstlisting}[language=C++,caption={等待条件写成谓词}]
this->ready_.wait(lock, [this] {
    return this->closed_ || !this->queue_.empty();
});
\end{lstlisting}

这行代码是协议文档。它告诉读者：消费者醒来的合法原因有两个，要么有任务，要么队列关闭。醒来后还要分支：

\begin{lstlisting}[language=C++,caption={醒来后区分任务和关闭}]
if (this->queue_.empty()) {
    return std::nullopt;
}
const int value = this->queue_.front();
this->queue_.pop();
return value;
\end{lstlisting}

如果忘了把 \code{closed_} 放进谓词，关闭空队列时消费者可能永远睡下去。如果忘了在 \code{close} 里 \code{notify_all}，等待者也可能看不到关闭事件。条件变量要同时检查三件事：状态在哪里改，谓词怎么写，谁负责通知。

\section{线程生命周期要有所有者}

启动线程也属于资源管理。裸 \code{std::thread} 如果对象析构时仍然 joinable，会调用 \code{std::terminate}。所以线程所有者必须定义退出协议：

\begin{lstlisting}[language=C++,caption={工作线程所有者的基本形状}]
class Worker {
public:
    Worker()
        : thread_{[this](std::stop_token token) {
              this->run(token);
          }}
    {
    }

private:
    void run(std::stop_token token)
    {
        while (!token.stop_requested()) {
            do_one_unit_of_work();
        }
    }

    std::jthread thread_;
};
\end{lstlisting}

\code{std::jthread} 帮你做析构时请求停止和 join，但它不能替你设计 \code{run} 如何从阻塞等待中醒来。如果工作线程在条件变量上等待，析构或关闭函数必须修改共享状态并通知。线程生命周期同样要有账本：谁启动，谁请求停止，谁唤醒，谁 join。

\section{本章验收：给每个共享状态标注保护者}

审查并发代码时，在每个共享字段旁边写一句“由谁保护”。例如：\code{queue_} 和 \code{closed_} 由 \code{mutex_} 保护；\code{requests_} 是 relaxed atomic 统计计数，不保护其他数据；\code{stop_} 是停止信号，但阻塞等待还需要条件变量通知。标注不出来的共享状态，就是事故候选。
